\section{The First Query Across Two Subgraphs}
\label{sec:ch09-the-first-query-across-two}

\index{seam!the first query that crosses one|(}%
\code{graph.yaml} gains an entry, pointing at the schema the new service
exports rather than at a copy of it:

\begin{minted}[highlightlines={7,8,9,10}]{yaml}
version: 1
subgraphs:
  - name: sessions
    routing_url: http://localhost:5001/graphql
    schema:
      file: ../src/Sessions/schema.graphql
  - name: speakers
    routing_url: http://localhost:5002/graphql
    schema:
      file: ../src/Speakers/schema.graphql
\end{minted}

\index{composition!the collision is gone}%
Compose it. The error section~\ref{sec:ch09-the-collision-that-is-still-there}
opened on is not there, and nothing replaced it:

\begin{minted}[fontsize=\footnotesize,breakanywhere]{text}
Router execution config successfully written to "router.json".
\end{minted}

\index{composition!silence is the whole result}%
That is the entire output. Two services that both name \code{Speaker}, and no
complaint, because only one of them claims to be able to produce one.

\subsection{The Query}

\index{seam!the query that crosses it}%
Start both services and the router, and send the request
chapter~\ref{ch:data-without-n-plus-1} has been sending since it was one
service:

\begin{minted}{text}
POST /graphql HTTP/1.1
Host: localhost:3002
Content-Type: application/json

{"query":"{ sessions(first: 10) { nodes { title speaker { name } } } }"}
\end{minted}

\begin{minted}[fontsize=\footnotesize]{json}
{"data":{"sessions":{"nodes":[
  {"title":"Schemas That Outlive Their Authors",
   "speaker":{"name":"Ada Fischer"}},
  {"title":"Reading a Query Plan",
   "speaker":{"name":"Ada Fischer"}},
  {"title":"Paging Without Offsets",
   "speaker":{"name":"Bruno Kaminski"}},
  {"title":"The Bank That Split Its Graph",
   "speaker":{"name":"Chidi Okafor"}}]}}}
\end{minted}

\index{seam!the response is byte for byte the old one}%
That is the answer chapter~\ref{ch:federation-model} printed, out of a graph
that did not exist then. Four sessions, four speaker names, each attached to
the right session, and nothing in the bytes saying that two processes and two
databases were involved. No identifier a client could use to redo the join for
itself, which was chapter~\ref{ch:federation-model}'s point about the response
and is still true of it.

\index{statement count!the seam does not add one}%
Now the number. Both services are running
chapter~\ref{ch:data-without-n-plus-1}'s statement log, so the question is
answerable rather than arguable:

\begin{center}
\begin{tabular}{lr}
\toprule
Service & Statements \\
\midrule
Sessions, on 5001 & 1 \\
Speakers, on 5002 & 1 \\
\midrule
Total & 2 \\
\bottomrule
\end{tabular}
\end{center}

\index{statement count!chapter 3's two, across a seam}%
Two. Chapter~\ref{ch:data-without-n-plus-1} measured this exact request at five
statements with a naive resolver and two with a DataLoader, inside one process,
and settled two as the baseline chapter~\ref{ch:router-n-plus-1} would measure
against. It is still two. What changed is that it is one in each of two
processes, and the second one is this:

\begin{minted}{text}
-- statement 1
SELECT "s"."Id", "s"."Bio", "s"."Name"
FROM "Speakers" AS "s"
WHERE "s"."Id" IN (@ids1, @ids2, @ids3)
\end{minted}

\index{DataLoader!the batch survives the move}%
Three parameters for four sessions, because two of them are Ada Fischer's.
That is chapter~\ref{ch:data-without-n-plus-1}'s batched lookup, in a different
process, against a different database, reached through a reference resolver
rather than through a field. The DataLoader did not stop working when the code
moved, which is worth saying out loud because the thing on the other side of it
changed completely: it used to be fed by one resolver running four times, and
it is now fed by one \code{\_entities} call carrying four representations.

\subsection{Reading the Plan}

\index{query plan!grows a second step}%
Chapter~\ref{ch:enter-the-router} said that every plan over a graph of one
subgraph is a \code{Sequence} holding one \code{Single}, and that the moment a
second name appeared in the subgraph list the plan is where you go to find out
what the router decided. Ask for it:

\begin{minted}[fontsize=\footnotesize]{json}
{
  "queryPlan": {
    "version": "1",
    "kind": "Sequence",
    "children": [
      {
        "kind": "Single",
        "fetch": {
          "kind": "Single",
          "subgraphName": "sessions",
          "subgraphId": "0",
          "fetchId": 0,
          "query": "{\n    sessions {\n        nodes {\n            title\n            speaker {\n                __typename\n                id\n            }\n        }\n    }\n}"
        }
      },
      {
        "kind": "Single",
        "fetch": {
          "kind": "BatchEntity",
          "path": "sessions.nodes.@.speaker",
          "subgraphName": "speakers",
          "subgraphId": "1",
          "fetchId": 1,
          "dependsOnFetchIds": [0],
          "representations": [
            {
              "kind": "@key",
              "typeName": "Speaker",
              "fragment": "fragment Key on Speaker {\n    __typename\n    id\n}"
            }
          ],
          "query": "query($representations: [_Any!]!){\n    _entities(representations: $representations){\n        ... on Speaker {\n            __typename\n            name\n        }\n    }\n}"
        }
      }
    ],
    "normalizedQuery": "{sessions {nodes {title speaker {name}}}}"
  }
}
\end{minted}

\index{query plan!the first fetch asks for a key, not a name}%
Read the first fetch before the second, because it already contains the
decision. The client asked for \code{speaker \{ name \}}. Sessions is asked for
\code{speaker \{ \_\_typename id \}}. The router rewrote the selection, because
it knows that service cannot answer \code{name} and knows what it will need in
order to ask somebody who can.

\index{BatchEntity@\code{BatchEntity}!the second word in the vocabulary}%
The second fetch is a \code{BatchEntity}, which is the seam in one word. Four
things in it carry the whole mechanism between them:

\begin{itemize}
  \item \code{path} is \code{sessions.nodes.@.speaker}. That is the field being
        filled in, with \code{@} standing for every element of the list. The
        router is not fetching speakers; it is fetching the contents of one
        field position in a response it has already partly built.
  \item \code{dependsOnFetchIds} is \code{[0]}. This step cannot start until
        the first has come back, because the keys it needs are in that answer.
        Nothing about the seam is parallel.
  \item \code{representations} is the shape of what will be sent, written as a
        fragment: a type name and a key, and nothing else.
  \item \code{query} is the \code{\_entities} call
        chapter~\ref{ch:first-subgraph} wrote a resolver for and could only
        exercise by hand. This is the caller it was waiting for.
\end{itemize}

\index{seam!costs one round trip, not one per row}%
The word that matters in \code{BatchEntity} is the first one. Four sessions
produce one fetch carrying four representations, not four fetches, and that is
the router's doing rather than anything the subgraph asked for.
Figure~\ref{fig:ch09-two-fetches} is the two shapes beside each other.

\begin{figure}[htbp]
  \centering
  % The plan grows a second step: chapter 8 answered one query with one fetch;
% here a second subgraph joins, and the same shape of client request now
% needs two fetches, the second of which cannot start until the first has
% come back.
%
% Same idiom as figures/tikz/ch01-one-field.tex and ch08-one-fetch.tex: each
% lane is a timeline read left to right, ticks mark stages, and a brace under
% the stages that matter carries a one-line label. Axis, tick, event, span
% and spanlabel styles are copied from those files unchanged.
%
% Lane A (sessions { nodes { title } }, answered inside one subgraph) and
% lane B (the same field plus speaker { name }, the first query this book
% answers across two subgraphs) share the same first four stages: client
% posts to the router, the router plans, the router posts to Sessions, and
% Sessions runs one statement. Those four ticks sit at the same x in both
% lanes for exactly that reason, so the eye reads the two rows as the same
% request up to the point where they part company. From there lane A is
% finished, so its axis simply stops; lane B keeps going for two more ticks,
% because that is the whole argument of the figure, and a lane that stopped
% at the same width as lane A would be hiding it. This is a deliberate
% widening, not a timing claim: SPEC decision 19 puts no milliseconds in
% this book, the horizontal axis is sequence only, and no gap here is sized
% to suggest one stage took longer than another.
%
% The dashed vertical rule is the seam idiom fixed by
% figures/tikz/ch05-the-join-moves.tex, reused unchanged: a dashed vertical
% marks where work crosses from one subgraph's process into another's. It
% appears once in lane B, between the statement that finishes the Sessions
% fetch and the tick where the router sends ids onward as representations,
% because that is the one crossing this request makes. Lane A crosses no
% seam and carries no dashed rule at all, which is half the figure's point:
% one lane has a seam to cross, the other does not.
%
% Two braces sit under lane B, with a visible gap on either side of the
% seam so neither brace is mistaken for spanning it, the same discipline
% ch08-one-fetch.tex uses for its own pair of braces. The shorter brace, under
% the fetch to Sessions, notes that the router rewrote the selection to ask
% for the id it will need rather than the name the client asked for. The
% main brace, under the fetch to Speakers, names the step lane A never
% takes at all: a fetch that depends on the first one returning, and that
% costs one more statement rather than one per session.
%
% No quotation marks anywhere in this file. The book's strict Ascii mode
% (see check-chapter.psd1) makes the prose gate read a literal " as a quote
% character, while \" is LaTeX's diaeresis accent and fails to compile
% inside a node. Every identifier below is written bare, and URL paths are
% broken onto their own line by hand rather than quoted or wrapped.
%
% Bare tikzpicture (house style): the figure environment, caption and label
% live at the call site in the section file.
\begin{tikzpicture}[
  x=1mm, y=1mm,
  every node/.append style={font=\sffamily\scriptsize},
  axis/.style={draw, line width=0.4pt,
    -{Straight Barb[length=1.4mm, width=1.6mm]}},
  tick/.style={draw, line width=0.4pt},
  event/.style={anchor=south, align=center, text width=20mm,
    inner sep=1pt},
  span/.style={draw, line width=0.4pt},
  spanlabel/.style={anchor=north, align=center, inner sep=2pt},
  lane/.style={anchor=west, font=\sffamily\scriptsize\itshape},
  boundary/.style={draw, line width=0.4pt, dashed},
]

% ------------------------------------------------------------- lane A
% At y=20, level with lane B's own label. The three-line event at x=36
% reaches y=13, so the 13 that ch08-one-fetch.tex uses for its lane A
% would sit on top of it here; that lane's tall labels start further
% right and its title has ended by then.
\node[lane] at (0,20) {A. sessions only: one query, one subgraph};

\draw[axis] (0,0) -- (90,0);
\foreach \x in {14,36,58,82}{
  \draw[tick] (\x,-1.5) -- (\x,1.5);
}

\node[event, text width=20mm] at (14,3)
  {client posts to\\localhost:3002\\/graphql};
\node[event, text width=22mm] at (36,3)
  {router plans:\\one fetch,\\subgraph sessions};
\node[event, text width=22mm] at (58,3)
  {router posts to\\localhost:5001\\/graphql};
\node[event, text width=18mm] at (82,3)
  {1 statement\\to SQLite};

% ------------------------------------------------------------- lane B
% Lane A carries no braces, so it needs less room beneath it than lane B
% does. 35mm puts the two lanes the same distance apart as the two in
% ch08-one-fetch.tex measured from axis to axis.
\begin{scope}[shift={(0,-35)}]
  \node[lane] at (0,20)
    {B. speaker name added: the first query split across two subgraphs};

  \draw[axis] (0,0) -- (150,0);
  \foreach \x in {14,36,58,82,112,140}{
    \draw[tick] (\x,-1.5) -- (\x,1.5);
  }

  \draw[boundary] (97,-2) -- (97,8);
  \node[spanlabel, inner sep=1pt] at (97,-2) {seam};

  \node[event, text width=20mm] at (14,3)
    {client posts to\\localhost:3002\\/graphql};
  \node[event, text width=20mm] at (36,3)
    {router plans:\\Sequence,\\two steps};
  \node[event, text width=22mm] at (58,3)
    {router posts to\\localhost:5001\\/graphql:\\id, not name};
  \node[event, text width=24mm] at (82,3)
    {1 statement\\to SQLite,\\Sessions service};
  \node[event, text width=30mm] at (112,3)
    {router sends ids\\as representations\\to localhost:5002\\/graphql};
  \node[event, text width=24mm] at (140,3)
    {1 statement\\to SQLite,\\Speakers service};

  \draw[span] (58,-4) -- (58,-7) -- (82,-7) -- (82,-4);
  \node[spanlabel] at (70,-8)
    {the router rewrote the selection:\\it asks Sessions for the id it\\will
     need, not the name asked for};

  \draw[span] (112,-4) -- (112,-7) -- (140,-7) -- (140,-4);
  \node[spanlabel] at (126,-8)
    {the step lane A never takes:\\a BatchEntity fetch that waits\\on the
     first fetch, and costs one\\more statement, not one per session};
\end{scope}

\end{tikzpicture}

  \caption{The same request before and after it has to cross a seam. Lane A
  stays inside one subgraph and the router adds nothing but a plan, which is
  where chapter~\ref{ch:enter-the-router} left it. Lane B is one field wider
  and costs a second fetch that cannot begin until the first returns, and the
  dashed rule is where the work leaves one service for another. Both lanes end
  in one statement per service, which is the part a DataLoader is responsible
  for and chapter~\ref{ch:router-n-plus-1} is about.}
  \label{fig:ch09-two-fetches}
\end{figure}

\index{router!batches representations, and the resolver may not}%
Hold on to the batching, because it is doing more work than it looks like.
The router sends one request per subgraph per step no matter how many rows are
in the list, so the count on the wire is fixed. What is not fixed is what the
service does with a list of four representations after it arrives, and
\code{SpeakerType}'s answer, one statement, is a property of the DataLoader
behind that resolver rather than of federation. Take the loader out and the
router's behavior does not change at all while the statement count does.
That is chapter~\ref{ch:router-n-plus-1}, and it is the first thing the next
part of the book measures.

\subsection{When One Of Them Is Not There}

\index{partial failure!one subgraph down is not the graph down}%
Chapter~\ref{ch:enter-the-router} stopped the only service there was and got a
null and an error. With three services the more useful question is what happens
when one of them goes and the others do not. Stop Speakers and send the same
request:

\begin{minted}[fontsize=\footnotesize]{json}
{"errors":[{"message":"Failed to fetch from Subgraph 'speakers' at Path 'sessions.nodes.@.speaker'."}],
 "data":{"sessions":{"nodes":[
   {"title":"Schemas That Outlive Their Authors","speaker":null},
   {"title":"Reading a Query Plan","speaker":null},
   {"title":"Paging Without Offsets","speaker":null},
   {"title":"The Bank That Split Its Graph","speaker":null}]}}}
\end{minted}

\index{partial failure!the error carries the path as well as the name}%
HTTP 200, every title present, every speaker null, and one error naming both
the subgraph that failed and the path it failed at. That path is new. When
there was one subgraph the error could only be about the whole request; with a
seam it can say which field position the hole is in, which is the difference
between a page that says \enquote{the graph is down} and a page that can render
the schedule and gray out the names.

\index{partial failure!a query that does not cross is untouched}%
Ask for something that does not cross the broken seam and there is no error at
all. Titles, abstracts and start times come back complete while Speakers is
still stopped, because every one of them routes to a service that is running.
The blast radius of a dead subgraph is the set of fields that route to it, and
no larger, provided those fields are nullable.
Chapter~\ref{ch:null-blast-radius} is what happens when they are not.
\index{seam!the first query that crosses one|)}%
